% 6.电路基础.tex 子文件，book文档类，主文档通过\input{6.电路基础.tex}导入
\chapter{第一章 电路元件与电路基本定律}

\section{电流、电压与电功率}
\subsection{电流}
荷电质点的有序运动形成电流。为了衡量电流的大小，在时间$dt$内，通过某截面的电荷量的代数和为$dq$，则电流定义为
\[
i=\frac{dq}{dt}
\]
其方向规定为正电荷运动的方向。

电荷的SI(国际单位制)单位为库[仑](符号为$\mathrm{C}$)，时间的SI单位为秒时，电流的SI单位为安[培](符号为$\mathrm{A}$)，$1\,\mathrm{A}=1\,\mathrm{C}/1\,\mathrm{s}$。

电流的方向既可用箭头”$\rightarrow$“表示，也可用双下标表示。例如在图1.1中电流可从a点流向b点或相反，分别用$i_{ab}$和$i_{ba}$表示。在分析复杂电路时，难以判断电流的真实方向，而且有些电路中电流的真实方向随时间不断改变。为此，在列写电路方程时，可任意假设一个电流的方向，称为电流的参考方向，并用箭头标示在图中，如图1.2所示。选定参考方向后，电流便是代数量。电流的真实方向由电流代数量的符号和参考方向来判断。当按参考方向求得的电流为正值时，表明其真实方向与参考方向一致；否则相反。

量值和方向不随时间变化的电流称为直流。

\subsection{电压}
单位正电荷从电路中的一点移动到另一点，电场力对电荷做功。设电场力，则电压定义为
\[
u_{ab}=\frac{dW_{ab}}{dq}
\]
在式中，$u_{ab}$表示a、b两点间的电压。

电压又称电位差，与电压相关的另一个物理量是电位。电位的物理意义是：在电路中可任选一点为参考点(零电位点,常用符号$\perp$表示)，电路中某点与参考点之间的电压称为该点的电位，用$\phi$表示。

电压的方向规定为从高电位指向低电位的方向。电压参考方向可用双下标表示。

\subsection{电动势}
从能量的角度看。设在$dt$的时间内，一电源使正电荷$dq$从负极经电源内部移至正极所做的功$dW$，电源的电动势定义为
\[
e=\frac{dW}{dq}
\]
即电源电动势。

\subsection{电功率}
电功率是用以衡量，则电功率定义为
\[
p=\frac{dW}{dt}
\]

当电能。

设某部分。

如果所选定的电压和电流的参考方向相反。

功率等于能量的转换速率，根据式可得在$t_0$时间内，电路吸收（电压、电流为关联参考方向时）或发出的能量为
\[
W(t)=\int_{t_0}^t p(\xi)\mathrm{d}\xi=\int_{t_0}^t u(\xi)i(\xi)\mathrm{d}\xi
\]
与判断功率。

发电站，若超过额定值，设备或部件。

\section{电阻元件}
在电路模型中，用电阻元件近似表示。在电路理论中，将流过相同电流的两个端子称为一个端口。线性二端电阻的端口电压与电流之间遵循欧姆定律，即二者成正比关系。

欧姆定律表达式为
\[
u=Ri \qquad \text{或} \qquad i=Gu
\]

若$u、i$参考方向相反。

电阻与电导是反映同一电阻元件性能而互为倒数的两个参数，即
\[
RG=1
\]
线性电阻的端口特性。

如果电阻$R=0$，则通过任意有限电流时其电压恒为零，故相当于短路，即短路线的电阻为零。

输入电阻的功率是$p=ui$，将式分别代入，得
\[
p=ui=Ri^2=Gu^2
\]

如果$R、G$是正值，则$p$非负，这样的电阻消耗电功率，属于耗能元件。从能量角度看，直到t瞬间输入电阻中的全部电能将是
\[
W=\int_{-\infty}^t p(\xi)\mathrm{d}\xi=\int_{-\infty}^t u(\xi)i(\xi)\mathrm{d}\xi=R
\]
此式就是物理学中的焦耳定律。导体。

对于非负的电阻，在任一时刻，式的结果非负。这表明：从全过程来看，这种元件只能吸收电能而不能发出电能。具有此种性质的元件，称为无源元件。

在电路模型中还存在一种负电阻。

从广义上讲，只要端口电压、电流关系是代数方程，就属于电阻元件，其电压、电流方程一般形式为
\[
u=u(i) \qquad i=i(u)
\]
上式中$u(i)$和$i(u)$分别表示。如果电压、电流不是正比关系，即电压、电流关系不是，就称为非线性电阻。例如，半导体二极管就属于非线性电阻，其电压、电流关系为
\[
i=I_0\left(e^{u/U_T}-1\right).
\]

\section{电容元件}
当在电容器两端通过电源施加电压时，两极板上分别聚集起等量的正、负电荷，并在极板之间建立电场，储存有电场能量。当电源移去后，电荷和电场可以继续存在，但不能无限期存在。电容能够储存电荷和电场能，是储能元件。从物理学已知当电容器填充线性介质（即介电常数与电场强度无关，如空气、橡皮、瓷、纱、胶木和云母等）时，它储存的电荷$q$与极板间电压$u$成正比。若电压的参考方向是从正极板指向负极板，则线性电容的$q-u$关系为
\[
q=Cu
\]
式中$C$是电容[系数]，由式可见其单位应是$\mathrm{C/V}$（库/伏），称为法[拉](符号为$\mathrm{F}$)。设极板上存在电荷$q$，再计算极板间电压$u$，二者之比即为电容。平行板电容器，其电容为
\[
C=\frac{\varepsilon A}{d}
\]
式中，$A$为每个极板面积，$d$为极板间距离，$\varepsilon$为介质的介电常数。对线性介质，电容$C$是常量，称为线性电容；否则称为非线性电容。

当$u$的参考方向是从正极板指向负极板时，则电流等于极板电荷的时间变化率
\[
i=\frac{dq}{dt}
\]
将，得
\[
i=C\frac{du}{dt}
\]
即线性电容的端口电流与端口电压的时间变化率成正比。当在端口上施加直流电压时，端口电流为零，等效成开路。

反过来，电荷$q$等于电流$i$对时间的积分，即
\[
q(t)=\int_{-\infty}^t i(\xi)\mathrm{d}\xi
\]
式表明$t$时刻电容上的电荷量是此时刻以前由电流充电和放电而“积累”的结果。或者说。所以电容具有记忆的性质，属于记忆元件。上式积分下限取$-\infty$，意指从很久以前算起，电容上的电荷为零。

由式，得
\[
u(t)=\frac{q(t)}{C}=\frac{1}{C}\int_{-\infty}^t i(\xi)\mathrm{d}\xi
\]
式即为线性电容的$u-i$关系，式中$u、i$应为关联参考方向；否则要在式中添一负号。

如果所研究
\[
u(t)=\frac{1}{C}\int_{-\infty}^{t_0}i(\xi)\mathrm{d}\xi+\frac{1}{C}\int_{t_0}^t i(\xi)\mathrm{d}\xi
=u(t_0)+\frac{1}{C}\int_{t_0}^t i(\xi)\mathrm{d}\xi
\]
式可见，电压$u(t_0)$表示$t_0$时刻以前全部电流对电压的影响，称为电容在$t_0$时的初始电压。

输入线性电容的功率为
\[
p=ui=Cu\frac{du}{dt}
\]
根据能量与功率的关系，由式可以求得截止到$t$时刻电容吸收的总能量$W_e(t)$。
在$t=-\infty$时电容尚未充电，所以$u(-\infty)=0$。故电容吸收的总能量为
\[
W_e(t)=\frac{1}{2}Cu^2=\frac{q^2}{2C}
\]
式就是电容储存电场能量的计算公式。

\section{电感元件}
在电机。尽管实际，但其共性是，在其周围激发磁场。

如果线圈的磁场存在于线性介质中，则磁链与电流成正比，称为线性电感；否则称为非线性电感。

设线性电感的磁链与电流的参考方向符合右手螺旋定则，则二者之间的关系为
\[
\Psi=Li
\]
式中$L$为电感[系数]。当磁链，磁链的单位是韦[伯](符号$\mathrm{Wb}$)，电流的单位是安[培]($\mathrm{A}$)时，电感的SI单位是。

根据电磁感应定律，磁链随时间的变化要在电感中产生感应电动势，记作$e$。当电压$u$、电动势$e$与电流$i$取关联参考方向，并且电流与
\[
u=-e=-\left(-\frac{d\Psi}{dt}\right)=\frac{d\Psi}{dt}
\]
对线性电感，将，得
\[
u=L\frac{di}{dt}
\]
即线性电感的端口电压与端口电流的时间变化率成正比。当电感中流过直流电流时，端口电压为零，等效成短路。

如果从电压求磁链或电流，则需积分，即
\[
\Psi(t)=\int_{-\infty}^t u(\xi)\mathrm{d}\xi
=\Psi(t_0)+\int_{t_0}^t u(\xi)\mathrm{d}\xi
\]
\[
i(t)=\frac{1}{L}\int_{-\infty}^t u(\xi)\mathrm{d}\xi
=i(t_0)+\frac{1}{L}\int_{t_0}^t u(\xi)\mathrm{d}\xi
\]
此两式表明，电感中某一瞬间的磁链和电流决定于此瞬间以前的全部电压，因此电感也属于记忆元件。

给出了线性电感上$u-i$的关系，它们应为关联参考方向；否则要在式中添一负号。

线性电感吸收的功率为
\[
p=ui=Li\frac{di}{dt}
\]
截止到$t$时刻电感吸收的能量为
\[
W_m = \int_{-\infty}^{i} p(\xi) \, d\xi
= L \int_{i(-\infty)}^{i} i(\xi) \, di(\xi)
= \frac{1}{2} L i^2 \bigg|_{i(-\infty)}^{i(t)}
\]
在$t=-\infty$时，电感电流为零，所以电感吸收能量为
\[
W_m=\frac{1}{2}Li^2=\frac{\Psi^2}{2L}
\]
式也是电感储存磁场能量的计算公式。当电感电流的绝对值增大时，所以电感元件也是一个储能元件。

\section{独立电源}
电源模型有电压源和电流源两种。

\subsection{电压源}
当电压源的电流在一定范围内时，它们提供的电压与电流大小无关，或为常量或随时间按确定规律变化。

电压源能够提供确定的电源电压$u_s$，称为源电压。所谓“确定”是指源电压$u_s$与流过电压源的电流无关，电压源的电流将由与其相连的电路来确定。$u_s$可以是常量，称为直流电压源（或恒定电压源），记作$u_s=U_s$；也可以是时变量，记作$u_s=u_s(t)$。源电压$u_s$的方向是从“$+$”级到“$-$”级。

电压源的端口特性可以用解析式表示
\[
u=u_s
\]
式中不含电流$i$，表示源电压电流的制约。源电压为常量时，端口电压与电流之间的关系可用所示的平行于电流坐标轴的直线表示，称为电源的端口特性曲线。当源电压$u_s=0$即电压源置零时，电压源的作用相当于短路。在图所示参考方向下，电压与电流形成的功率是
\[
p=u_si
\]
由于这里$u_s、i$取非关联参考方向，在电源内部电流。

\subsection{电流源}
电流源能够提供确定的端口电流$i_s$，称为源电流。这里“确定”是指$i_s$与电流源端口电压无关，电流源的端口电压决定于它所接的外电路。$i_s$可以是常量即恒定电流，记作$i_s=I_s$，也可以是时变量，记作$i_s=i_s(t)$。

电流源端口特性的解析表达式为
\[
i=i_s
\]
式不含电压$u$，表明电流源的端口电流不受端口电压的制约。如果源电流$i_s$为常量，则式对应$u-i$平面上一条平行于电压坐标轴的的直线。当$i_s=0$即电流源置零时，电流源的作用相当于开路。

在图所示参考方向下，电流源输出的功率为
\[
p=ui_s
\]
当特性曲线位于第$\mathrm{I}、$第$\mathrm{III}$象限时，上式$p>0$，电流源处于供电状态；当特性曲线位于第$\mathrm{II}、$第$\mathrm{IV}$象限时，上式$p<0$，电流源处于充电状态。

电压源和电流源作为元件模型。

\section{受控电源}
在电路模型中还要用到另一类电源，它们的源电压或源电流受电路中另一处的电压或电流控制，这类电源称为受控电源或非独立电源。

受控源是一种用来表示。本书只讨论线性受控源，即源电压（流）与控制电压（源）成正比关系。用菱形表示受控源，以便与独立源相区分。

受控源的端口方程为
\begin{itemize}
    \item 电压控制电压源
    \item 电流控制电压源
    \item 电压控制电流源
    \item 电流控制电流源
\end{itemize}
式中带脚注$S$的变量表示源，带脚注$C$的变量表示控制量。参量$\alpha,r,g,\beta$都是常量。

\section{基尔霍夫定律}
1847年基尔霍夫建立了电路结构对电流、电压的约束关系。分别称为基尔霍夫电流定律和基尔霍夫电压定律，统称基尔霍夫定律。

\subsection{电路结构}
电路元件通过端子相互连接便形成具有特定结构的电路。电路结构。

本书将每个二端元件称为一条支路，若干支路的连接点称为节点，其中两个支路的连接点称为简单节点。

\subsection{基尔霍夫电流定律}
基尔霍夫电流定律表述为：在集中参数电路中，任一时刻流出（或流入）任一节点的支路电流代数和等于零，即
\[
\sum i_k=0
\]
通常规定当式中$i_k$参考方向为流出节点时。

式称为节点的KCL方程，它反映了与任意一个节点相连的所有支路电流所满足的关系。

KCL还可推广用于流出（或流入）闭合边界的支路电流，表述为：在集中参数电路中，任一时刻流出（或流入）任一闭合边界$S$的支路电流代数和等于零，即
\[
\]
式称为广义基尔霍夫电流定律。

如将流出。

这表明：任一时刻，流出任一节点（或闭合边界）电流的代数和等于流入该节点电流的代数和，即
\[
\]

\subsection{基尔霍夫电压定律}
基尔霍夫电压定律表述为：在集中参数电路中，任一时刻沿任一回路各支路电压的代数和等于零，即
\[
\]
通常规定当式中$u_k$的参考方向与回路方向相同时，$u_k$前面。

式称为回路的KVL方程，它是任一回路所包含的支路电压必须满足的约束。

推广到一般情况就是：沿任一回路，各支路电压降的代数和等于电压升的代数和，即
\[
\]
再推广到一般情况就是：集中参数电路中,任意两点之间的电压具有确定值，与计算路径无关。

对电路中的每个网孔或回路均可列出KVL方程，但所有这些。可以验证：任一回路的设电路有$b$个支路$n$个节点，可以证明：平面电路的网孔数即独立KVL方程的个数等于$b-(n-1)$。

基尔霍夫定律与集中参数电路元件性质无关。

\endinput